\documentclass[10pt,a4paper]{amsart}
\usepackage[margin=19mm]{geometry}
\usepackage{amsmath,amssymb,mathtools}
\usepackage[scr=rsfs,cal=euler]{mathalfa}
\usepackage{xfrac}
\usepackage[unicode,hidelinks,bookmarks=false]{hyperref}
\newcommand{\ie}{i.e.\ }
\newcommand{\cf}{cf.\ }
\newcommand{\resp}{respectively\ }
\newcommand{\Subsection}{\subsection}
\providecommand{\nobreakdash}{\nobreak\hbox{-}}
\setlength{\emergencystretch}{2em}
\multlinegap0pt
\usepackage{amssymb}
\usepackage{xcolor}
\usepackage{nicefrac}
\usepackage{caption}
\usepackage{natbib}
\usepackage{cancel}
\usepackage{bbm}
\usepackage{multirow}
\usepackage{siunitx}
\usepackage{algorithm}\usepackage{algpseudocode}
\usepackage[shortlabels]{enumitem}
\newtheorem{proposition}{Proposition}
\newtheorem{lemma}{Lemma}
\newcommand{\E}{\mathbb{E}}
\newcommand{\R}{\mathbb{R}}
\def\aGmin{a_{\mathcal{G}}}
\def\aStarmax{a_{+}^\star}
\def\aStarmin{a_{-}^\star}
\def\amax{a_{+}}
\def\bmax{b_T^{+}}
\def\cGmax{c_{\mathcal{G}}^{+}}
\def\cGmin{c_{\mathcal{G}}}
\def\cStarmax{c_{+}^\star}
\def\cStarmin{c_{-}^\star}
\def\cTmax{c_T^{+}}
\def\cmax{c_{+}}
\title{Schr\"odinger bridge problem via empirical risk minimization}
\author{Denis Belomestny, Alexey Naumov, Nikita Puchkin, Denis Suchkov}
\date{Source: arXiv:2602.08374v1 (CC BY 4.0)}
\begin{document}
\maketitle

\noindent\emph{Source and licence.} Schr\"odinger bridge problem via empirical risk minimization. Version arXiv:2602.08374v1 (CC BY 4.0), distributed by arXiv under the Creative Commons Attribution 4.0 International licence, http://creativecommons.org/licenses/by/4.0/, as recorded in the arXiv licence metadata for this exact version and shown on the abstract page https://arxiv.org/abs/2602.08374v1. Full licence notice: \texttt{licenses/arxiv-2602.08374-CC-BY-4.0.txt}.\par\smallskip

\noindent\textit{Editorial scope.} Counted unit: the article's lemma properties (source0.tex lines 1196--1221). The article's complete original proof of it is delivered with it (source0.tex lines 1222--1329, one \texttt{proof} environment of 108 lines). No mathematics is written, repaired or re-derived: the statement and the proof are the article's own lines, in the article's own notation, and no retained line is substituted or re-flowed.  Retained uncounted as the source-local context the proof needs: lines 905--924; lines 955--964; lines 997--1004; lines 1057--1064.  Nothing was omitted from any retained span, so the package carries no omission record.  No retained line contains a citation, so the wrapper carries no bibliography and no bibliography entry.  No retained line contains a figure environment, an image-inclusion command or drawing code, so no image is delivered and the package declares no asset dependency.  Editorial formatting only, in addition: an AMS \texttt{amsart} preamble that restates the article's own declarations and macros used by the retained lines, namely the environments \texttt{proposition}, \texttt{lemma} on separate counters, together with the standard packages those lines need.  The retained environments print in this document as Proposition 1, Lemma 1 in order of appearance, whereas the article prints the same items with the numbering of its own full document.  The complete native source of the article as distributed in the arXiv e-print for this exact version is bundled unchanged as \texttt{sources/194241/source0.tex}.
\begin{proposition}[Two--sided Gaussian bounds for the fixed point]
\label{prop:g-two-sided-subgaussian-direct}
Assume \emph{(Q)}, \emph{(R0)}, \emph{(RT)}. Let $g^\star:\mathbb{R}^d\to(0,\infty)$ be measurable and satisfy
\begin{equation}\label{eq:fixed-point-normalization}
  g^\star(y) = \mathcal{C}[g^\star](y), \qquad y\in\mathbb{R}^d,
  \qquad\text{and}\qquad
  \int_{\mathbb{R}^d} \frac{1}{g^\star(z)}\,\rho_T(z)\,dz = 1.
\end{equation}
Set \(\aStarmin := 2a_-,\) \(\aStarmax := \frac{a_+}{2}.\)
Then there exist constants $\cStarmin, \cStarmax>0$ depending only on
$q_T,\rho_0,\rho_T$ (but not on $g^\star$), such that
\begin{equation}\label{eq:g-two-sided-Gaussian}
  \cStarmin \exp(-\aStarmin \|y\|^2)
  \;\le\;
  g^\star(y)
  \;\le\;
  \cStarmax \exp(-\aStarmax \|y\|^2),
  \qquad \forall y\in\mathbb{R}^d.
\end{equation}
\end{proposition}

\begin{equation}
\label{D upper boound}
  D_{g^\star}(x)
  =
  \int q_T(x,z)\,\frac{\rho_T(z)}{g^\star(z)}\,dz
  \le
  \cmax \int \frac{\rho_T(z)}{g^\star(z)}\,dz
  =
  \cmax.
\end{equation}

\begin{equation}
\label{gauss density formula}
\int_{\mathbb{R}^d}
  \exp(-\gamma\|z\|^2)\,dz
  =
  \left(\frac{\pi}{\gamma}\right)^{d/2}
  <\infty,
\end{equation}

\begin{equation}
     \label{Dmin def}
     D_{g^\star}(x)
  =
  m(x)\,\mathbb{E}_{\nu_x}\Bigl[\frac{1}{g^\star(Z)}\Bigr]
  \ge
  m_-\,k_- =: \underline{D}^\star > 0
\end{equation}

\begin{lemma}
\label{lemm: key properties}
Suppose that the assumptions \emph{(Q)}, \emph{(R0)}, \emph{(RT)} and \emph{(G)} hold.  Assume additionally that $\bmax> 2 \aGmin$. Then there exist constants $\underline{D}, \overline{D}, L_{D,\infty}, L_{\mathcal C, \infty} > 0$ such that
\begin{align}
\label{D bounds}
  & \underline{D}
  \;\le\;
  D_g(x)
  \;\le\;
  \overline{D},
  \qquad x\in B(x_0,R_0),\; g\in\mathcal{G}; \\
  \label{D Lip}
  & \|D_g-D_h\|_{L^\infty(B(x_0, R_0))}
  \le
  L_{D,\infty} \|g-h\|_{\infty};  \\
  \label{C lip infty}
  & \|\mathcal C[g]-\mathcal C[h]\|_\infty
  \le
  L_{\mathcal C, \infty}\,\|g-h\|_{L^\infty(\R^d)}; \\
  \label{C infty}
  & \|\mathcal C[g]\|_\infty
  \le
  (\cmax/\underline{D})\,\|g\|_{L^\infty(\R^d)}\, .
\end{align}
Note that $\underline{D}, \overline{D}$ do not depend on the class $\mathcal G$. 
\end{lemma}

\begin{proof}
The proof of \eqref{D bounds} repeats the proof of Proposition \ref{prop:g-two-sided-subgaussian-direct}. We make necessary changes. First, we replace \eqref{D upper boound} by
\begin{align*}
  D_{g}(x)
  &=
  \int q_T(x,z)\,\frac{\rho_T(z)}{g(z)}\,dz
  \le
  \cmax \cTmax \cGmin^{-1}\int 
  \exp\!\left(-\Bigl(\bmax- \aGmin\Bigr)\|z\|^2\right)\,dz
  \\
  & =
  \cmax \cTmax  \cGmin^{-1} \left(\frac{\pi}{\bmax- \aGmin}\right)^{d/2} = : \overline{D}\, .
\end{align*}
Recall the proof of \eqref{Dmin def}. Note that
$$
\E_{\rho_T}[1/g] \geq (\cGmax)^{-1}\,, 
$$
and we can take $\underline{D}: = m_-\,k_- (\cGmax)^{-1}$. 
We prove \eqref{D Lip}. By definition,
\[
  D_g(x)-D_h(x)
  =
  \int_{\mathbb{R}^d} q_T(x,z)\,\rho_T(z)\,
  \Bigl(\frac{1}{g(z)}-\frac{1}{h(z)}\Bigr)\,dz.
\]
Using
\[
  \left|\frac{1}{g(z)}-\frac{1}{h(z)}\right|
  =
  \frac{|g(z)-h(z)|}{g(z)h(z)}
  \le
  \|g-h\|_\infty\cdot \frac{1}{g(z)h(z)}
\]
and the lower bound in \emph{(G)} (applied to both $g$ and $h$),
we obtain for all $z$,
\[
  \frac{1}{g(z)h(z)}
  \le
  \cGmin^{-2}\exp(2 \aGmin \|z\|^2).
\]
Hence
\begin{equation}\label{eq:Dg-Dh-basic}
  |D_g(x)-D_h(x)|
  \le
  \|g-h\|_\infty\,
  (\cGmin)^{-2}
  \int_{\mathbb{R}^d}
    q_T(x,z)\,\rho_T(z)\,\exp(2 \aGmin \|z\|^2) dz.
\end{equation}
We now bound the integral uniformly in $x\in B(x_0, R_0)$. By \emph{(Q)} and the inequality
\[
  \|x-z\|^2 \ge \frac{1}{2}\|z\|^2 - \|x\|^2,
\]
we have, for all $x\in B(x_0, R_0)$,
\[
  q_T(x,z)
  \le
  \cmax \exp(-\amax\|x-z\|^2)
  \le \cmax \exp(\amax R_\star^2)\exp\!\left(-\frac{\amax}{2}\|z\|^2\right)
\]
with $R_\star:=\|x_0\|+R_0.$
Combining this with \emph{(RT)} yields
\[
  q_T(x,z)\,\rho_T(z)\,\exp(2 \aGmin \|z\|^2)
  \le
  \cmax
  \cTmax \exp(\amax R_\star^2)\,
  \exp\!\left(-\Bigl(\bmax- 2 \aGmin +\frac{\amax}{2}\Bigr)\|z\|^2\right).
\]
By \eqref{gauss density formula}, uniformly for $x\in B(x_0, R_0)$,
\[
  |D_g(x)-D_h(x)|
  \le
  L_{D,\infty} \|g-h\|_\infty \, ,
\]
where 
\begin{equation}
\label{LD definition}
L_{D,\infty}
  :=
  \cGmin^{-2}
  \cmax
  \cTmax \exp(\amax R_\star^2) 
  \left(\frac{\pi}{\bmax- 2 \aGmin +\amax/2}\right)^{d/2}\, .
\end{equation}
This proves \eqref{D Lip}. For each $y \in \R^d$,
\[
  \mathcal C[g](y)-\mathcal C[h](y)
  =
  \int_{\mathrm{supp}(\rho_0)}\rho_0(x)\,q_T(x,y)
  \left(\frac{1}{D_g(x)}-\frac{1}{D_h(x)}\right)\,dx.
\]
By \eqref{D bounds} and \eqref{D Lip} we get 
\begin{align*}
  |\mathcal C[g](y)-\mathcal C[h](y)|
  & \le
  \frac{L_{D,\infty} \|g-h\|_\infty}{\underline D^2}
  \int_{\mathrm{supp}(\rho_0)}\rho_0(x)\,q_T(x,y)\,dx \\
  & \le L_{\mathcal C, \infty} \|g-h\|_\infty, 
\end{align*}
where 
\begin{eqnarray}
\label{CD definition}
    L_{\mathcal C, \infty} = L_{D,\infty} \cmax \, .
\end{eqnarray}
This proves \eqref{C lip infty}. 

\end{proof}

\end{document}
